\section{Efficience des ressources et numérique responsable}

Au-delà de la sécurisation du flux, la réécriture des synchronisations participe à une utilisation plus efficiente des ressources. Les fichiers déjà présents ne sont plus recopiés systématiquement : leur taille et leur date de modification sont vérifiées avant tout transfert. Certaines productions sont également écrites dans des tampons locaux avant une livraison réseau contrôlée. Pour des fichiers raster parfois volumineux, cette organisation limite les transferts, les écritures sur les stockages partagés et les traitements inutiles.

Je n'ai toutefois pas réalisé de mesure énergétique ni de bilan environnemental de la chaîne. Je présente donc cette démarche comme une optimisation concrète des ressources techniques, et non comme une évaluation complète de son impact écologique.
